\begin{frame}[plain]
    \begin{center}
        \LARGE Introdução
    \end{center}
\end{frame}
% \begin{frame}{Introdução}
%     Redes neurais convolucionais (\textit{convolutional neural network} ou \textit{CNN}) são uma combinação de redes neurais artificiais com o processo de convolução.
% \end{frame}
% \begin{frame}{Analogia ao neurônio e cérebro}
% \end{frame}
\begin{frame}[plain]
    \begin{figure}
        \centering
        \caption{Representação de uma ligação entre neurônios}
        \includegraphics[height=7cm]{images/introduction/neuron.png}
        \begin{center}
            {\footnotesize Fonte: BruceBlaus, 2013}
        \end{center}
        \label{fig:enter-label}
    \end{figure}
\end{frame}
% \begin{frame}{Redes neurais}
% \end{frame}
\begin{frame}[plain]
    \begin{figure}
        \centering
        \caption{Representação de uma rede neural}
        \includegraphics[width=\linewidth]{images/introduction/neural network.png}
        \begin{center}
            {\footnotesize Fonte: Pramoditha, 2022}
        \end{center}
        \label{fig:enter-label}
    \end{figure}
\end{frame}
% \begin{frame}{Redes neurais convolucionais}
% \end{frame}
\begin{frame}[plain]
    \begin{figure}
        \centering
        \caption{Representação de uma redes neural convolucional}
        \includegraphics[height=7cm]{images/introduction/cnn.jpg}
        \begin{center}
            {\footnotesize Fonte: Saha, 2018}
        \end{center}
        \label{fig:enter-label}
    \end{figure}
\end{frame}